% Rédaction : SOMBIE Roseline ; relecture : BANDE Hamidou
\chapter{Workflow KNIME de contrôle et de transformation}\label{chap:knime}

\section{Rôle dans la chaîne}
Le workflow KNIME répond à l'exigence E2 : contrôler et transformer les données brutes avant leur
chargement. Il lit \texttt{data/raw/transactions\_brutes.csv} (1010 lignes) et écrit deux fichiers
dans \texttt{data/processed/} :
\begin{itemize}
  \item \texttt{transactions\_valides.csv} : 946 lignes, plus la colonne \texttt{categorie\_montant} ;
  \item \texttt{transactions\_rejetees.csv} : 64 lignes, chacune avec son \texttt{motif\_rejet}.
\end{itemize}
KNIME écrit des CSV plutôt que d'écrire directement dans PostgreSQL. Aucun pilote JDBC n'est à
configurer sur le poste, et la sortie peut être vérifiée automatiquement \emph{avant} d'être chargée
(chapitre~\ref{chap:donnees}).

Le workflow a été construit avec KNIME Analytics Platform 5.12, sans extension, en suivant la
spécification \texttt{docs/KNIME\_WORKFLOW.md}. Il est exporté dans \texttt{knime/transactions\_etl.knwf}.

\capture{knime_01_workflow.png}{Le workflow complet : 23 nœuds exécutés (feux verts)}{fig:knime-workflow}

\section{Les étapes du workflow}
Les 23 nœuds forment six étapes. Chaque nœud porte son numéro dans son nom, pour que la capture se
lise dans l'ordre.

\begin{table}[H]
\centering\small
\caption{Les six étapes du workflow KNIME}
\label{tab:knime-etapes}
\begin{tabularx}{\linewidth}{|l|l|X|}
\hline
\rowcolor{bluedark}\textcolor{white}{\textbf{Nœuds}} & \textcolor{white}{\textbf{Étape}} & \textcolor{white}{\textbf{Ce qui est fait}} \\
\hline
1 & Lire & CSV Reader ; \textbf{toutes les colonnes forcées en texte} (String) \\
\hline
2--8 & Normaliser & codes en majuscules sans espaces ; montant : espaces retirés, virgule $\rightarrow$ point ; date française \texttt{JJ/MM/AAAA hh:mm:ss} $\rightarrow$ ISO \texttt{AAAA-MM-JJThh:mm:ss} \\
\hline
9--10 & Contrôler & Rule Engine : une colonne \texttt{motif\_rejet} ; séparation des lignes \texttt{OK} et des rejets \\
\hline
11--13 & Dédoublonner & doublons sur \texttt{reference} seule ; la première occurrence est gardée, les suivantes reçoivent \texttt{DOUBLON\_REFERENCE} \\
\hline
14--18 & Enrichir et écrire & montant converti en nombre ; \texttt{categorie\_montant} ; ordre des colonnes ; écriture des valides \\
\hline
19--23 & Écrire les rejets & réunion des deux sorties de rejet ; écriture ; comptage par motif \\
\hline
\end{tabularx}
\end{table}

\paragraph{Pourquoi tout lire en texte.} Si le CSV Reader devine les types, il lit \texttt{montant}
comme un nombre : les valeurs \texttt{"mille"} ou \texttt{"12500,00"} deviennent des valeurs
manquantes. On ne pourrait plus dire \emph{pourquoi} la ligne est fausse. Le montant n'est converti
en nombre qu'au nœud 15, une fois la ligne contrôlée.

\paragraph{Corriger ou rejeter.} Une anomalie dont l'intention est certaine est corrigée : minuscules,
espaces, virgule décimale, séparateur de milliers, date au format français. Les 26 anomalies de ce
type sont acceptées après correction. Une donnée absente, fausse ou ambiguë est rejetée avec un motif.

\section{Les règles de contrôle (nœud 9)}
Le nœud 9 est le cœur du workflow. Ses 30 règles sont évaluées dans l'ordre et la première règle
vraie donne le motif. Elles sont la traduction exacte de la fonction \texttt{motifRejet} de
\texttt{data/scripts/rules.js}, dans le même ordre : c'est ce qui permet de comparer KNIME à l'ETL
de référence.

\begin{lstlisting}[caption={Début des règles du nœud 9 (Rule Engine)}]
MISSING $reference$ => "REFERENCE_MANQUANTE"
$reference$ = "" => "REFERENCE_MANQUANTE"
NOT $reference$ MATCHES "TXN-[0-9]{6}" => "REFERENCE_INVALIDE"
...
$montant$ LIKE "-*" => "MONTANT_NEGATIF"
NOT $montant$ MATCHES "[0-9]+([.][0-9]{1,2})?" => "MONTANT_NON_NUMERIQUE"
...
$compte_destination$ = $compte_source$ => "COMPTES_IDENTIQUES"
TRUE => "OK"
\end{lstlisting}

\capture{knime_02_noeud9.png}{Configuration du nœud 9 : les 30 règles de contrôle}{fig:knime-noeud9}

\section{Résultat obtenu}
Les points de contrôle intermédiaires ont tous été atteints :

\begin{table}[H]
\centering\small
\caption{Points de contrôle du workflow (poste de SOMBIE, 29/10/2026)}
\label{tab:knime-controles}
\begin{tabular}{|l|l|r|}
\hline
\rowcolor{bluedark}\textcolor{white}{\textbf{Nœud}} & \textcolor{white}{\textbf{Contrôle}} & \textcolor{white}{\textbf{Lignes}} \\
\hline
1 CSV Reader & lignes lues, 10 colonnes en texte & 1010 \\
\hline
10 Séparation & sortie \texttt{OK} / rejets du nœud 9 & 956 / 54 \\
\hline
13 Séparation & valides / doublons & 946 / 10 \\
\hline
23 GroupBy & motifs différents (total 64) & 18 \\
\hline
\end{tabular}
\end{table}

\begin{table}[H]
\centering\small
\caption{Rejets par motif produits par KNIME}
\label{tab:knime-motifs}
\begin{tabular}{|l|r|l|r|}
\hline
\rowcolor{bluedark}\textcolor{white}{\textbf{Motif}} & \textcolor{white}{\textbf{Lignes}} & \textcolor{white}{\textbf{Motif}} & \textcolor{white}{\textbf{Lignes}} \\
\hline
DOUBLON\_REFERENCE & 10 & DATE\_INVALIDE & 4 \\
\hline
MONTANT\_NEGATIF & 5 & REFERENCE\_MANQUANTE & 4 \\
\hline
COMPTE\_SOURCE\_MANQUANT & 4 & TYPE\_INVALIDE & 4 \\
\hline
COMPTE\_SOURCE\_INVALIDE & 3 & DATE\_MANQUANTE & 3 \\
\hline
DESTINATION\_MANQUANTE & 3 & DEVISE\_INVALIDE & 3 \\
\hline
MONTANT\_MANQUANT & 3 & MONTANT\_NON\_NUMERIQUE & 3 \\
\hline
MONTANT\_NUL & 3 & REFERENCE\_INVALIDE & 3 \\
\hline
STATUT\_INVALIDE & 3 & CANAL\_INVALIDE & 2 \\
\hline
COMPTES\_IDENTIQUES & 2 & DESTINATION\_INVALIDE & 2 \\
\hline
\multicolumn{3}{|l|}{\textbf{Total des rejets} (54 anomalies + 10 doublons)} & \textbf{64} \\
\hline
\end{tabular}
\end{table}

\capture{knime_03_rejets_par_motif.png}{Sortie du nœud 23 : 18 motifs de rejet, 64 lignes au total}{fig:knime-motifs}

\section{Vérification automatique}
La commande \texttt{npm run data:compare} compare la sortie de KNIME, ligne par ligne et champ par
champ, au résultat de l'ETL de référence (\texttt{data/expected/}). Elle est obligatoire avant tout
chargement.

\begin{lstlisting}[caption={Vérification de la sortie KNIME}]
$ npm run data:compare
[comparaison] CONFORME : 946 valides, 64 rejetées, motifs identiques au résultat attendu.
\end{lstlisting}

\capture{knime_04_conforme.png}{La sortie de KNIME est conforme au résultat attendu}{fig:knime-conforme}

Pour rejouer la vérification sur un autre poste, on importe \texttt{knime/transactions\_etl.knwf}
dans KNIME, puis on adapte les chemins du CSV Reader (nœud 1) et des deux CSV Writer (nœuds 18 et 22).

\section{Difficultés rencontrées}
La construction a buté sur trois points. Aucun n'a demandé de modifier les règles : la règle du
projet est de corriger le workflow, jamais les règles, pour « faire passer » la comparaison.

\begin{table}[H]
\centering\small
\caption{Difficultés rencontrées lors de la construction du workflow}
\label{tab:knime-difficultes}
\begin{tabularx}{\linewidth}{|l|X|X|}
\hline
\rowcolor{bluedark}\textcolor{white}{\textbf{Nœud}} & \textcolor{white}{\textbf{Symptôme}} & \textcolor{white}{\textbf{Cause et correction}} \\
\hline
2 & 17 colonnes au lieu de 10 & le mode par défaut ajoute des colonnes suffixées \texttt{\_transformed} ; choisir \emph{Replace selected input columns} \\
\hline
12 & règle en rouge & dans KNIME 5.12, la colonne de statut s'appelle \texttt{Duplicate Status} et non \texttt{duplicate-type-classifier} comme prévu ; insérer la colonne depuis la liste plutôt que la taper \\
\hline
18 & le nœud refuse de s'exécuter & un port optionnel \emph{File System Connection} avait été ajouté par erreur et restait non relié ; le retirer \\
\hline
\end{tabularx}
\end{table}

La spécification \texttt{docs/KNIME\_WORKFLOW.md} reste la référence pour reconstruire le workflow
sur un autre poste. Son tableau « écart $\rightarrow$ cause $\rightarrow$ correction » couvre les erreurs que la comparaison
peut signaler.
